\chapter{Congenital disorder}
\index{Congenital disorder}

\section*{Definition and Overview}
A congenital disorder, also referred to as a congenital anomaly, birth defect, or congenital malformation, is a structural, functional, or metabolic abnormality that is present at birth. These disorders can range from minor physical anomalies with minimal clinical significance to severe structural defects or metabolic abnormalities that cause profound disability or mortality. Congenital disorders may be identified prenatally, at birth, or later in infancy or childhood as physiological systems mature and developmental milestones are assessed. They represent a major global public health challenge, acting as a leading cause of infant mortality, childhood morbidity, and long-term disability, thereby requiring coordinated multidisciplinary medical, surgical, and supportive care.

\section*{Epidemiology}
Globally, congenital disorders affect approximately 2\% to 3\% of all live births, resulting in millions of affected infants annually. In many countries, the prevalence is consistent with these figures, with around 3\% of neonates diagnosed with one or more congenital anomalies. Congenital heart defects are the most common structural disorders, followed by neural tube defects and chromosomal conditions such as Down syndrome (trisomy 21). The epidemiology of congenital disorders is influenced by maternal age (with older maternal age associated with higher rates of chromosomal aneuploidies), genetic factors, environmental exposures, maternal nutritional status, and access to comprehensive prenatal screening and care. In resource-limited settings, the burden is disproportionately high due to lack of preventive measures and delayed access to corrective surgeries.

\section*{Aetiology and Pathophysiology}
The developmental origins of congenital disorders are multifactorial, arising from a complex interplay of genetic factors, maternal health, environmental exposures, and nutritional influences during critical windows of embryonic and foetal development.

\begin{itemize}
    \item \textbf{Genetic and chromosomal abnormalities:} Monogenic mutations, microdeletions, or chromosomal aneuploidies that alter cellular signaling, gene expression, and organogenesis during embryonic development.
    \item \textbf{Teratogenic exposures:} Gestational exposure to harmful substances, including certain medications (such as sodium valproate or thalidomide), alcohol (causing foetal alcohol spectrum disorder), tobacco, recreational drugs, and environmental chemicals.
    \item \textbf{Maternal infections:} Vertical transmission of pathogens during pregnancy, commonly grouped under the TORCH acronym (Toxoplasmosis, Other agents like syphilis and parvovirus B19, Rubella, Cytomegalovirus, and Herpes simplex virus), which directly damage developing foetal tissues.
    \item \textbf{Nutritional deficiencies:} Inadequate maternal intake of essential micronutrients, particularly folate deficiency, which impairs neural tube closure and increases the risk of spina bifida and anencephaly.
    \item \textbf{Maternal chronic illnesses:} Poorly controlled pre-existing maternal conditions, such as pregestational diabetes mellitus or maternal phenylketonuria, which alter the intrauterine environment and disrupt normal embryological pathways.
\end{itemize}

\section*{Clinical Features}
The clinical presentation of a congenital disorder is highly variable, depending on the specific organ systems affected, the severity of the structural or functional defect, and the developmental stage of the patient.

\begin{itemize}
    \item \textbf{Structural malformations:} Visible physical defects including cleft lip or cleft palate, limb abnormalities (such as polydactyly or syndactyly), abdominal wall defects (like gastroschisis or omphalocele), or open neural tube defects.
    \item \textbf{Cardiovascular signs:} Cyanosis (bluish skin discolouration), cardiac murmurs, tachypnoea, feeding difficulties, and poor weight gain in neonates, indicating underlying cardiac shunts or valvular abnormalities.
    \item \textbf{Neurological deficits:} Muscle hypotonia or hypertonia, microcephaly, hydrocephalus, abnormal reflexes, and delayed attainment of developmental milestones.
    \item \textbf{Sensory impairments:} Congenital cataracts, corneal clouding, microphthalmia, or failure to respond to visual and auditory stimuli, indicating congenital blindness or sensorineural hearing loss.
    \item \textbf{Metabolic and systemic signs:} Failure to thrive, persistent neonatal jaundice, vomiting, unusual body odour, or extreme lethargy, suggestive of inborn errors of metabolism.
\end{itemize}

\section*{Differential Diagnosis}
Diagnosing congenital disorders requires distinguishing them from acquired perinatal conditions, transient physiological changes, and late-onset genetic syndromes.

\begin{itemize}
    \item \textbf{Acquired perinatal injuries:} Hypoxic-ischaemic encephalopathy, birth trauma (such as neonatal clavicle fractures or brachial plexus injuries), and early-onset neonatal sepsis that manifest shortly after birth but do not originate from prenatal developmental anomalies.
    \item \textbf{Transient physiological adaptations:} Normal neonatal phenomena such as physiological jaundice, transient tachypnoea of the newborn (TTN), or mild developmental dysplasia of the hip that corrects spontaneously.
    \item \textbf{Late-onset genetic disorders:} Hereditary neurological, metabolic, or muscular disorders (such as Huntington's disease or certain dystrophies) that are genetically predetermined but present with no clinical features at birth.
    \item \textbf{Postnatal environmental and nutritional delays:} Developmental delays, cognitive deficits, or failure to thrive caused by postnatal neglect, lack of environmental stimulation, or severe nutritional deprivation rather than congenital anomalies.
\end{itemize}

\section*{When to Seek Medical Care}
Parents, caregivers, and primary health providers should monitor neonates and infants closely for signs of developmental delay, feeding difficulties, or structural abnormalities. Timely evaluation by a paediatrician is indicated if any congenital disorder is suspected.

\textbf{Contact your local emergency services for:}
\begin{itemize}
    \item Severe respiratory distress, including grunting, nasal flaring, chest wall retractions, or prolonged pauses in breathing (apnoea).
    \item Central cyanosis, characterised by a bluish discolouration of the lips, tongue, and trunk.
    \item Extreme lethargy, difficulty waking the infant for feeding, or a floppy, unresponsive muscle tone.
    \item Seizures, rhythmic jerking of the limbs, or abnormal eye deviations.
    \item Forceful, projectile, or green (bile-stained) vomiting.
\end{itemize}

\section*{Diagnosis and Investigations}
A definitive diagnosis is achieved through a combination of prenatal screening, detailed physical examination, biochemical tests, and advanced imaging modalities.

\begin{itemize}
    \item \textbf{Prenatal screening tests:} Non-invasive prenatal testing (NIPT) measuring cell-free foetal DNA, maternal serum screening, and high-resolution obstetric ultrasound to identify structural abnormalities or soft markers of chromosomal disorders.
    \item \textbf{Invasive prenatal diagnostic procedures:} Chorionic villus sampling (CVS) or amniocentesis to collect foetal cells for karyotyping, fluorescence in situ hybridisation (FISH), or chromosomal microarray analysis.
    \item \textbf{Newborn screening programme:} Blood spot screening (heel prick test) performed 48 to 72 hours after birth to screen for cystic fibrosis, congenital hypothyroidism, phenylketonuria, and other metabolic disorders.
    \item \textbf{Postnatal imaging:} Echocardiography for suspected cardiac defects, cranial ultrasound or magnetic resonance imaging (MRI) for neurological anomalies, and skeletal surveys for bone dysplasias.
    \item \textbf{Clinical genetic testing:} Whole exome or genome sequencing to identify specific pathogenic genetic variants in infants presenting with complex, syndromic, or multiple congenital anomalies.
\end{itemize}

\section*{Management}
Management plans must be highly individualised, requiring a multidisciplinary team including paediatricians, surgeons, geneticists, allied health professionals, and support networks.

\begin{itemize}
    \item \textbf{Surgical intervention:} Corrective, reconstructive, or palliative surgeries, such as closure of congenital heart defects, cleft lip and palate repair, spinal closure for spina bifida, or orthopaedic correction of talipes equinovarus.
    \item \textbf{Pharmacological and medical therapies:} Hormone replacement therapy (e.g., levothyroxine for congenital hypothyroidism), enzyme replacement therapies, or specific dietary formulations to bypass blocked metabolic pathways.
    \item \textbf{Allied health and rehabilitation:} Early intervention services incorporating physiotherapy, occupational therapy, and speech pathology to optimise physical mobility, communication, and cognitive milestones.
    \item \textbf{Assistive technologies:} Fitting of orthotics, prosthetics, hearing aids, cochlear implants, or visual aids to enhance functional independence.
    \item \textbf{Genetic counselling:} Educating families regarding the recurrence risk of specific congenital disorders in future pregnancies and discussing reproductive options.
\end{itemize}

\section*{Complications}
Congenital disorders can lead to various long-term physical, mental, and social complications.

\begin{itemize}
    \item \textbf{Chronic physical disability:} Long-term impairment of mobility, coordination, or sensory function, limiting daily activities and independent living.
    \item \textbf{Intellectual and cognitive impairment:} Developmental delays and learning difficulties that impact educational attainment and adaptive behaviours.
    \item \textbf{Recurrent infections:} Increased susceptibility to respiratory tract infections (common in congenital heart disease and cystic fibrosis) or urinary tract infections (due to renal tract anomalies).
    \item \textbf{Organ failure:} Progressive deterioration of vital organs, such as congestive heart failure, renal failure, or respiratory insufficiency, resulting from uncorrected structural anomalies.
    \item \textbf{Nutritional deficiencies and failure to thrive:} Chronic feeding difficulties, gastrointestinal malabsorption, and growth restriction.
    \item \textbf{Caregiver burnout and psychosocial distress:} Psychological impact, including anxiety, depression, and social isolation among affected individuals, parents, and families.
\end{itemize}

\section*{Prognosis and Outcomes}
The prognosis for individuals with congenital disorders varies widely depending on the organ systems involved, the severity of the anomalies, the presence of multiple defects, and the timeliness of clinical intervention. While some severe chromosomal and structural abnormalities are life-limiting and may lead to stillbirth or early neonatal death, major advances in prenatal diagnostics, neonatal intensive care, and paediatric surgery have significantly improved survival rates and long-term outcomes. Many children with mild-to-moderate congenital disorders can achieve a normal life expectancy and enjoy a high quality of life, particularly when supported by early intervention programmes, educational assistance, and continuous, coordinated medical care.

\section*{Prevention and Self-Care}
While not all congenital disorders can be prevented, preconception and antenatal health strategies can significantly reduce the risk of occurrence.

\begin{itemize}
    \item \textbf{Preconception folic acid supplementation:} Daily consumption of 400 to 500 micrograms of folic acid starting at least one month prior to conception and continuing through the first trimester to prevent neural tube defects.
    \item \textbf{Preconception vaccination:} Ensuring immunity to vaccine-preventable infections, particularly rubella, varicella, and influenza, before becoming pregnant.
    \item \textbf{Avoidance of teratogenic substances:} Abstaining from alcohol, tobacco, recreational drugs, and non-essential prescription medications during planning and throughout pregnancy.
    \item \textbf{Optimal maternal disease control:} Managing pre-existing chronic conditions, such as diabetes mellitus, hypertension, and epilepsy, under close medical supervision before and during pregnancy.
    \item \textbf{Safe environmental practices:} Avoiding exposure to hazardous chemicals, pesticides, heavy metals, and unnecessary radiation during gestation.
\end{itemize}

\section*{Sources and Further Reading}
The following reputable organisations provide further information and clinical resources on congenital disorders:

\begin{itemize}
    \item World Health Organization. \textit{Fact sheet on congenital anomalies, global statistics, and prevention strategies.} https://www.who.int/
    \item Healthdirect Australia. \textit{Comprehensive guide to congenital anomalies, screening, and local support services.} https://www.healthdirect.gov.au/
    \item Better Health Channel. \textit{Information on birth defects, genetic counselling, and maternal health guidelines.} https://www.betterhealth.vic.gov.au/
    \item Centers for Disease Control and Prevention. \textit{Research, surveillance, and prevention recommendations for birth defects.} https://www.cdc.gov/
    \item National Health and Medical Research Council. \textit{Australian clinical guidelines on preconception care and pregnancy advice.} https://www.nhmrc.gov.au/
\end{itemize}
